\documentclass[reqno,10pt,a4paper]{amsart}
\usepackage{mathtools,amssymb,microtype}
\usepackage[T1]{fontenc}
\usepackage{lmodern}
\usepackage[a4paper,textwidth=15.7cm,hcentering,top=1.8cm,bottom=2.1cm,headheight=10pt,headsep=14pt]{geometry}
\usepackage{xcolor}
\definecolor{linkblue}{RGB}{0,61,105}
\usepackage[colorlinks=true,linkcolor=linkblue,citecolor=linkblue,urlcolor=linkblue]{hyperref}
\setlength{\emergencystretch}{1em}
\numberwithin{equation}{section}
\newcommand{\R}{\mathbb R}
\newcommand{\T}{\mathbb T}
\linespread{1.075}
\title[Long-time behavior and stability]{Long-time behavior and stability of\\ kinetic and dispersive equations}
\author{Wenrui Huang}
\date{}
\hypersetup{pdftitle={Project Description -- Long-time behavior and stability of kinetic and dispersive equations},pdfauthor={Wenrui Huang}}
\begin{document}
\maketitle
\begin{center}
{\small Herchel Smith Research Fellowship in Pure Mathematics \enspace$\cdot$\enspace Project Description}
\end{center}
\smallskip
\noindent\textbf{Research areas.} Nonlinear partial differential equations, kinetic theory, and dispersive equations.
\medskip

My research concerns kinetic and dispersive partial differential equations. A central example in my work is the Vlasov--Poisson system, which models collisionless plasmas in the repulsive case and self-gravitating systems in the attractive case. I study the long-time behavior of solutions and the stability of stationary states under small perturbations. I propose to study scattering under external potentials, the effect of specular reflection on the asymptotic dynamics of a plasma, and damping and nonlinear stability near small inhomogeneous equilibria.

\section{Long-time behavior and scattering for particles and waves}

\subsection*{Riesz interactions and the Coulomb threshold.}
I study how the interaction potential $\Gamma$ and the external potential $V$ shape the asymptotic dynamics of collisionless kinetic equations in Hamiltonian form,
\[
 \partial_t\mu+\{\mu,\mathcal H\}=0,\qquad
 \mathcal H=\tfrac12|v|^2+V(x)-\lambda\phi,\qquad
 \phi=\Gamma*\rho,\qquad \rho=\int_{\R^d}\mu^2\,dv .
\]
Here $\{\cdot,\cdot\}$ is the Poisson bracket in $(x,v)$ and $\mu^2$ is the phase-space density. For the Riesz interaction $\Gamma(x)=|x|^{-\alpha}$ in three dimensions with $V=0$, we compare the solution with free transport. Along dispersing trajectories the force is of order $t^{-1-\alpha}$. The velocity converges, but its correction, of order $t^{-\alpha}$, accumulates in the position. The deviation from free motion is bounded for $\alpha>1$, of order $\log t$ at the Coulomb exponent $\alpha=1$, and of order $t^{1-\alpha}$ for $\alpha<1$. In joint work with Hyunwoo Kwon~\cite{HK24}, I obtained scattering for small data and constructed wave operators for Riesz interactions near the Coulomb exponent. A kinetic pseudoconformal transformation sends infinite time to a finite endpoint, where a symplectic change of variables controls the singularity. This also solves the inverse problem of constructing a solution with a prescribed asymptotic profile.

\subsection*{The Vlasov--Hartree system and a common structure for modified scattering.}
In joint work with Mengyi Xie~\cite{HX26a,HX26b}, I studied a mixed classical--quantum model of a Bose--Einstein condensate interacting with a degenerate Fermi gas. The particle density generates the potential acting on the wave, and the wave density generates the force on the particles. The trajectory correction and the wave phase must therefore be constructed together. For small, regular, localized data and $1/2<\alpha<2$, we identified three scattering regimes for both components. The coupling is integrable along free rays for $\alpha>1$. At $\alpha=1$ the particles acquire a logarithmic displacement and the wave a logarithmic phase. For $1/2<\alpha<1$, a time-dependent symplectic transformation removes the leading trajectory correction and a regularized effective phase controls the wave.

For $0<\alpha\leq1/2$, we treat the Hartree equation, the Vlasov--Riesz equation, and their coupling in a common framework. Along rays $x=ty$, a nonlinear action $\theta_k(t,y)$ corrects a kinetic profile by a momentum translation $\nabla_y\theta_k$ and a wave profile by a phase modulation. The momentum shift itself produces a quadratic term of order $t^{-2\alpha}$, whose time integral diverges in this range. At the same order, the action changes the density and hence the potential. With $m_k$ the rescaled density of component $k$, the action and density satisfy
\[
 \partial_t\theta_k-\frac{|\nabla_y\theta_k|^2}{2t^2}=t^{-\alpha}\varphi_j,
 \qquad
 \partial_t m_k-\frac1{t^2}\nabla_y\!\cdot(m_k\nabla_y\theta_k)
 =-\frac1{t^2}\nabla_y\!\cdot\mathcal J_k .
\]
Here $t^{-\alpha}\varphi_j$ is the regularized potential generated by the source density $m_j$, and $\mathcal J_k$ is a residual current. In a single equation $j=k$, while in the coupled system the other component supplies the source. The density drives the action, whose gradient carries the leading density flux. Iterating the two equations produces the trajectory and phase corrections $t^{1-n\alpha}$ for $n\alpha<1$, with a logarithm when $n\alpha=1$. Each coefficient is determined recursively from the limiting densities.

\subsection*{External potentials.}
For Vlasov--Poisson with a repulsive harmonic potential $V(x)=-|x|^2/2$, Kwon and I introduced a hyperbolic lens transform and constructed wave operators and scattering maps~\cite{HK26}. The change of time $s=\tanh t$ sends $t=\infty$ to $s=1$, absorbs the external potential, and preserves the Hamiltonian structure. The transformed force carries the factor $(1-s^2)^{(d-4)/2}$. It is integrable up to $s=1$ for $d\geq3$, which gives linear scattering, and borderline nonintegrable in $d=2$, which produces a logarithmic correction in the lens variables. In one dimension I introduced a Hamilton--Jacobi correction that shifts the momentum by the gradient of an action forced by the self-consistent potential~\cite{H26}. For repulsive self-interactions this gives modified scattering for data of arbitrary size. In a separate dispersive problem I constructed modified wave operators for the Hartree equation with a repulsive external Coulomb potential and small, localized scattering data~\cite{H24}.

\emph{Proposed research.} The corresponding dynamics for an attractive harmonic potential remain largely unexplored. Our lens transform also applies in the attractive setting and may provide a useful starting point. The unperturbed orbits are periodic, and I would like to understand how the self-consistent force accumulates over successive periods. A broader goal is to identify the radial external potentials for which solutions exhibit scattering or modified scattering. For radially symmetric solutions, the motion under the external Hamiltonian at fixed angular momentum is integrable, and the asymptotic action--angle variables developed in~\cite{PW21,PWY22} may be useful. The problem is substantially harder for nonradial solutions, even when the external potential remains radial, because the self-consistent force can also change angular momentum. In this setting, the Coulomb and harmonic potentials are distinguished by their superintegrable Hamiltonian flows.

\section{Asymptotic stability of the vacuum in a convex domain}

The vacuum, in which no particles are present, is the simplest plasma equilibrium. Its asymptotic stability asks whether a sufficiently dilute plasma remains globally regular and approaches a simpler transport regime at large times. With Benoit Pausader and Masahiro Suzuki~\cite{HPS24}, I studied this question for the three-dimensional repulsive Vlasov--Poisson system in a class of unbounded convex domains $D$. The wall is perfectly conducting, so the electrostatic potential satisfies a Dirichlet boundary condition. Particles are absorbed when they reach the wall. The presence of the wall breaks translation invariance, and the transport boundary is characteristic along grazing trajectories.

A central geometric object in our analysis is the asymptotic domain. After placing the origin on the boundary, we define
\[
 D_\infty=\operatorname{Int}\bigcap_{s>0}s^{-1}D .
\]
This cone describes the velocity directions compatible with escape to infinity. For sufficiently small, regular, localized data satisfying the boundary compatibility conditions, we obtain a unique global classical solution. The potential and electric field decay at the optimal rates $t^{-1}$ and $t^{-2}$, respectively. The mass carried by velocities outside $\overline{D_\infty}$ tends to zero. If $D_\infty$ is empty, the total surviving mass tends to zero. Otherwise the particles with limiting velocities in $D_\infty$ exhibit logarithmically modified scattering. Their limiting distribution and the Dirichlet Green function of $D_\infty$ together determine the leading electric field.

The proof combines convex geometry, elliptic theory, and kinetic transport. We extend the Dirichlet potential across the wall and use a grazing function to quantify proximity to tangential impact. In free-transport coordinates, we combine weighted phase-space estimates with a scale decomposition of the Green function. The rescaled domains $t^{-1}D$ then identify the leading field along an escaping ray as $E(tv,t)\sim t^{-2}E_\infty(v)$. Its tail produces a $t^{-1}$ velocity correction and hence the logarithmic displacement of the surviving particles.

\emph{Proposed research.} A natural next step is to study a perfectly reflecting boundary, for which particles undergo specular reflection. With unit normal $n(x)$, the velocity after reflection is $v-2(v\cdot n(x))n(x)$. In the absorbing problem a characteristic ends at its first impact. With reflection, a near-tangential trajectory can return many times, and the field generated between impacts changes the angle of the next collision. A first question is whether the grazing estimates can be propagated across successive reflections with sufficient control of derivatives of the characteristic flow. I want to understand how repeated interactions with the boundary affect global regularity, asymptotic stability, and the long-time distribution of particles.

\section{Stability and instability of small BGK waves}

Bernstein--Greene--Kruskal (BGK) waves are spatially inhomogeneous steady states of the Vlasov--Poisson system. In the one-dimensional periodic problem, an even steady state has the form
\[
 \mu_S(x,v)=F(\mathcal H_S(x,v)),\qquad
 \mathcal H_S(x,v)=\tfrac12v^2-\phi_S(x).
\]
After normalizing $\min\phi_S=0$, the energy level $\mathcal H_S=0$ is the separatrix between trapped particles and passing particles. Varying the energy distribution $F$ gives a local family of stationary states with infinitely many tangent directions, even after fixing the period and position of the well.

With Dongfen Bian, Emmanuel Grenier, and Benoit Pausader~\cite{BGHP26}, I studied small one-well BGK waves bifurcating near a marginal homogeneous equilibrium $\mu_0$. For perturbations with the period of the wave, the spectral criterion is given by the dielectric function of $\mu_0$. Under the generic one-well hypotheses, including $\mu_0''(0)\ne0$ and the absence of nonzero neutral frequencies, a zero of the dielectric function in the open lower half-plane yields an exponentially growing linear mode. In the absence of such a zero, the wave is linearly stable up to infinite-dimensional modulation. We use action--angle coordinates and prove quantitative resolvent estimates for the dielectric operator. The expansion of this operator exhibits a transition at frequencies of order $\sqrt{\varepsilon}$, where $\varepsilon$ is the amplitude of the wave and trapped particles can compete with passing ones.

On the stable branch, admissible zero-charge perturbations decompose into a variation within the family of stationary states, symmetry terms, and a dynamical part. A perturbation can change the energy distribution or give the wave a drifting center. Passing particles can also carry an odd current without generating an electric field. After isolating these contributions, the dynamical electric field satisfies
\[
 \int_0^\infty\|E^{\mathrm{dyn}}(t)\|_{L_x^2}^2\,dt<\infty .
\]

\emph{Proposed research.} For the linearly stable waves, we are working toward quantitative damping estimates and, under suitable regularity assumptions, exponential decay of the electric field. A longer-term goal is to establish nonlinear stability. Here modulation must account for changes in the energy distribution as well as the position and velocity of the wave. In the unstable regime, the next step is to promote the exponentially growing linear mode to a nonlinear instability result, using ideas developed by Guo and Strauss and by Lin~\cite{GS95,L05}. A further question concerns the degenerate case $\mu_0''(0)=0$, which is excluded by the generic hypotheses. Understanding this regime may require a higher-order expansion of the Hamiltonian and of the associated dispersion operator.

At Cambridge, I would particularly welcome the opportunity to discuss these questions with Cl\'ement Mouhot, whose work on Landau damping is closely related to the stability problems above, and to interact with Pierre Rapha\"el on nonlinear dispersive equations.

\begin{thebibliography}{99}
\setlength{\itemsep}{1pt}\setlength{\parskip}{0pt}\setlength{\parsep}{0pt}
\bibitem{HK24} W. Huang and H. Kwon. \href{https://doi.org/10.1007/s00023-026-01735-7}{Scattering of the Vlasov--Riesz System in the Three Dimensions}. \emph{Ann. Henri Poincar\'e} (2026).
\bibitem{HX26a} W. Huang and M. Xie. Long-time dynamics of Vlasov--Hartree systems across the Coulomb threshold. \href{https://arxiv.org/abs/2608.04402}{arXiv:2608.04402}.
\bibitem{HX26b} W. Huang and M. Xie. A common structure for modified scattering: Vlasov--Riesz, Hartree, and their coupling. \href{https://arxiv.org/abs/2610.02643}{arXiv:2610.02643}.
\bibitem{HK26} W. Huang and H. Kwon. Scattering map for the Vlasov--Poisson system with a repulsive harmonic potential. \href{https://arxiv.org/abs/2602.21344}{arXiv:2602.21344}.
\bibitem{H26} W. Huang. Asymptotic behavior of the one-dimensional Vlasov--Poisson system with a repulsive harmonic potential. \href{https://arxiv.org/abs/2610.02551}{arXiv:2610.02551}.
\bibitem{H24} W. Huang. Modified wave operators for the Hartree equation with repulsive Coulomb potential. \href{https://arxiv.org/abs/2411.02737}{arXiv:2411.02737}.
\bibitem{HPS24} W. Huang, B. Pausader, and M. Suzuki. \href{https://doi.org/10.1007/s00220-026-05731-0}{The Vlasov--Poisson System with a Perfectly Conducting Wall: Convex Domains}. \emph{Comm. Math. Phys.} \textbf{407}, 219 (2026).
\bibitem{BGHP26} D. Bian, E. Grenier, W. Huang, and B. Pausader. Stability and instability of small BGK waves. \href{https://arxiv.org/abs/2601.10030}{arXiv:2601.10030}.

\bibitem{PW21} B. Pausader and K. Widmayer,
\href{https://link.springer.com/article/10.1007/s00220-021-04117-8}{\emph{Stability of a point charge for the Vlasov--Poisson system: The radial case}},
Comm. Math. Phys. \textbf{385} (2021), no.~3, 1741--1769.

\bibitem{PWY22} B. Pausader, K. Widmayer, and J. Yang,
\href{https://doi.org/10.4171/JEMS/1518}{\emph{Stability of a point charge for the repulsive Vlasov--Poisson system}},
J. Eur. Math. Soc. \textbf{28} (2026), 2751--2848.

\bibitem{GS95} Y. Guo and W.~A. Strauss,
\href{https://onlinelibrary.wiley.com/doi/abs/10.1002/cpa.3160480803}{\emph{Instability of BGK equilibria}},
Comm. Pure Appl. Math. \textbf{48} (1995), no.~8, 861--894.

\bibitem{L05} Z. Lin,
\href{https://onlinelibrary.wiley.com/doi/abs/10.1002/cpa.20028}{\emph{Nonlinear instability of periodic BGK waves for the Vlasov--Poisson system}},
Comm. Pure Appl. Math. \textbf{58} (2005), no.~4, 505--528.
\end{thebibliography}

\end{document}
